\documentclass{ltjsarticle}

%%%%%%%%%%packages%%%%%%%%%%
%% colors and links
\usepackage[svgnames]{xcolor}
\usepackage[colorlinks,citecolor=DarkGreen,linkcolor=DeepPink,linktocpage,unicode]{hyperref} 

%% equations
%%%% math
\usepackage{amsmath,amsfonts,amssymb,amsthm}
\usepackage{mathtools}
\usepackage{mathrsfs}
\usepackage{bm}
\usepackage{cancel}
\usepackage{dsfont}
%%%% physics
\usepackage{siunitx}
\usepackage{physics}
%%%% chemistry
\usepackage[version=3]{mhchem}

%% positioning
\usepackage{array}
\usepackage{float}

%% table
\usepackage{booktabs}
\usepackage{multirow}
\usepackage{hhline}
\usepackage{caption}
\captionsetup{format=hang}
\usepackage{subcaption}

%% figure
\usepackage{graphicx}
\usepackage{tikz}

%% decorations
\usepackage{titlesec}
\usepackage{picture}

%% framing
\usepackage{fancybox}
\usepackage{boites}
\usepackage{tcolorbox}
\tcbuselibrary{skins,theorems,breakable}

%% citation
\usepackage{cite}

%% miscellaneous
%%%% comment
\usepackage{comment}
%%%% Japanese
\usepackage{ascmac}
\usepackage{listings} %日本語のコメントアウトをする場合jlisting（もしくはjvlisting）が必要
\lstset{
    basicstyle={\ttfamily},
    identifierstyle={\small},
    commentstyle={\smallitshape},
    keywordstyle={\small\bfseries},
    ndkeywordstyle={\small},
    stringstyle={\small tfamily},
    frame={tb},
    breaklines=true,
    columns=[l]{fullflexible},
    numbers=left,
    xrightmargin=0pt,
    xleftmargin=27pt,
    numberstyle={\scriptsize},
    stepnumber=1,
    lineskip=-0.5ex
}

%%macros
\usepackage{mylogics}
\usepackage{myalgebra}
\usepackage{mygeometry}
\usepackage{myQM}

%%%%%%%%%%optional settings%%%%%%%%%%

%%%%%図表並列%%%%%
\makeatletter
\newcommand{\figcaption}[1]{\def\@captype{figure}\caption{#1}}
\newcommand{\tblcaption}[1]{\def\@captype{table}\caption{#1}}
\makeatother

%%%%%itemization%%%%%
\renewcommand{\labelenumi}{\theenumi}
\renewcommand{\theenumi}{(\arabic{enumi})} % 箇条書きをローマ数字に

%%%%%%theorem environments%%%%%
\newtheoremstyle{mystyle}%   % スタイル名
    {}%                      % 上部スペース
    {}%                      % 下部スペース
    {\normalfont}%          % 本文フォント
    {}%                      % インデント量
    {\bf}%                  % 見出しフォント
    {.}%                      % 見出し後の句読点
    {\newline}%                     % 見出し後のスペース
    {\underline{\thmname{#1}\thmnumber{#2}\thmnote{（#3）}}}%
    % 見出しの書式 (can be left empty, meaning `normal')
\theoremstyle{mystyle} % スタイルの適用

\newtheorem{theorem}{定理}[section]
\newtheorem{definition}{定義}[section]
\newtheorem{proposition}[definition]{命題}
\newtheorem{corollary}[theorem]{系}
\renewcommand{\proofname}{証明}

%proof
\makeatletter % use at mark
\renewcommand{\qedsymbol}{$\blacksquare$}
\renewenvironment{proof}[1][\proofname]{\par
    \pushQED{\qed}%
    \normalfont \topsep6\p@\@plus6\p@\relax
    \trivlist
    \item[\hskip\labelsep
        \itshape
    \textbf{\underline{#1}}]\ignorespaces
    % {\bf\underline{#1}\@addpunct{.}}]\ignorespaces % ピリオドあり
}{%
    \popQED\endtrivlist\@endpefalse
}
\makeatother % end at mark

%%%%%%mathtools%%%%%
\mathtoolsset{showonlyrefs=true} % 被参照数式のみ数式番号割り振り
\numberwithin{equation}{section}

\makeatletter
\@addtoreset{equation}{section}
\makeatother

%%%%%%framing%%%%%

%%%%%title decorations%%%%%

%% section
% \titleformat{\section}[block]
% {}{}{0pt}
% {
%     \colorbox{teal}{\begin{picture}(0,20)\end{picture}}
%     \hspace{0pt}
%     \normalfont \Huge\bfseries \thesection
%     \hspace{0.5em}
% }
% [
% \begin{picture}(100,0)
%     \put(3,30){\color{teal}\line(1,0){400}}
% \end{picture}
% \
% \vspace{-50pt}
% ]

% % subsection
% \titleformat{\subsection}[block]
% {}{}{0pt}
% {
%     \colorbox{blue}{\begin{picture}(0,10)\end{picture}}
%     \hspace{0pt}
%     \normalfont \Large\bfseries \thesubsection
%     \hspace{0.5em}
% }
% [
% \begin{picture}(100,0)
%     \put(3,18){\color{blue}\line(1,0){300}}
% \end{picture}
% \
% \vspace{-30pt}
% ]

% % subsubsection
% \titleformat{\subsubsection}[block]
% {}{}{0pt}
% {
%     \normalfont \large\bfseries \thesubsubsection
%     \hspace{0.5em}
% }
% [
% \begin{picture}(100,0)
%     \put(3,15){\color{black}\line(1,0){200}}
% \end{picture}
% \
% \vspace{-25pt}
% ]

% section
% \titleformat{\section}[block]
% {}{}{0pt}
% {
%     \colorbox{blue}{\begin{picture}(0,10)\end{picture}}
%     \hspace{0pt}
%     \normalfont \Large\bfseries \thesection
%     \hspace{0.5em}
% }
% [
% \begin{picture}(100,0)
%     \put(3,18){\color{blue}\line(1,0){300}}
% \end{picture}
% \\
% \vspace{-30pt}
% ]

% % subsection
% \titleformat{\subsection}[block]
% {}{}{0pt}
% {
%     \normalfont \large\bfseries \thesubsection
%     \hspace{0.5em}
% }
% [
% \begin{picture}(100,0)
%     \put(3,15){\color{black}\line(1,0){200}}
% \end{picture}
% \\
% \vspace{-25pt}
% ]
\mathchardef\mhyphen="2D

\newcommand{\lto}{\longrightarrow}
\newcommand{\lmto}{\longmapsto}
\newcommand{\btl}{\blacktriangleleft}
\newcommand{\btr}{\blacktriangleright}
\usepackage{tabularx}
\newcolumntype{Y}{>{\arraybackslash}X}

\newcommand{\TBA}{TBA}

\newcommand{\spkA}{安藤貴政}
\newcommand{\spkB}{緒方芳子}
\newcommand{\spkC}{菅野浩明}
\newcommand{\spkD}{前川拓海}
\newcommand{\spkE}{笠真生}
\newcommand{\spkF}{和田博貴}

\newcommand{\instA}{YITP}
\newcommand{\instB}{RIMS}
\newcommand{\instC}{名古屋大学}
\newcommand{\instD}{理化学研究所}
\newcommand{\instE}{プリンストン大学}
\newcommand{\instF}{東北大学}

\newcommand{\titleA}{Anomalous Symmetries on Topological States}
\newcommand{\titleB}{\TBA}
\newcommand{\titleC}{位相的弦理論と Gopakumar-Vafa 不変量}
\newcommand{\titleD}{Topics from higher topoi}
\newcommand{\titleE}{Multipartite entanglement and phases of matter}
\newcommand{\titleF}{Topological Aspects of Worldsheet Theories}

\newcommand{\abst}[5]{
    \Large
    \textbf{#1}
    \normalsize
    
    \vspace{10pt}

    \textbf{講演者} #2（#3）

    \textbf{講演時間} #4

    \vspace{5pt}

    % \textbf{概要} 
    \begin{quote}
        #5
    \end{quote}

    \vspace{10pt}
}


\begin{document}

\title{RFCMP2026-S}
\author{京都大学 基礎物理学研究所 湯川記念館1F パナソニック国際交流ホール}
\date{2026年5月12日 (火) - 15日 (金)}
\maketitle

本研究会はハイブリッド形式（対面 \& zoom）で実施します．参加登録は \url{https://indico.global/event/16468/} から行なってください\footnote{オンライン参加情報や限定公開の録画は\textbf{参加登録者のみ}に共有いたします．}．

% \section*{講演者（予定）}

% \begin{itemize}
%     \item \spkA（\instA）
%     \item \spkB（\instB）
%     \item \spkC（\instC）
%     \item \spkD（\instD）
%     \item \spkE（\instE）
%     \item \spkF（\instF）
% \end{itemize}


\section*{5月12日 (火)}
\vspace{-10pt}
\begin{table}[H]
    \centering
    \begin{tabularx}{0.8\linewidth}{llY}
        \toprule
        \multicolumn{1}{l}{\textbf{時間}}
        &\multicolumn{1}{l}{\textbf{講演者}}
        &\multicolumn{1}{l}{\textbf{タイトル}} \\
        \midrule
        10:00-13:00 & &受付・自由議論 \\
        13:00-14:30 &\spkE & \titleE \\ %\titleAA                                        \\
        15:00-16:30 &\spkE & \titleE \\ %\href{https://youtu.be/iSlKx8mIJzs}{\titleB}    \\
        17:00-18:30 &\spkF & \titleF \\ %\href{https://youtu.be/1ug-Wihe9xs}{\titleBB}   \\
    \end{tabularx}
\end{table}%

\section*{5月13日 (水)}
\vspace{-10pt}
\begin{table}[H]
    \centering
    \begin{tabularx}{0.8\linewidth}{llY}
        \toprule
        \multicolumn{1}{l}{\textbf{時間}}
        &\multicolumn{1}{l}{\textbf{講演者}}
        &\multicolumn{1}{l}{\textbf{タイトル}} \\
        \midrule
        10:00-11:30 &\spkF &\titleF   \\ %\href{https://youtu.be/K9ZWQPkibxk}{\titleC} \\
        11:30-13:00 & & 昼休憩 \\
        13:00-14:30 &\spkA &\titleA  \\ %\href{https://youtu.be/rhHHrmUZW7s}{\titleC} \\
        15:00-16:30 &\spkA &\titleA  \\ %\href{https://youtu.be/xow1kzEf_n4}{\titleD} \\
        16:30-18:00 & & 自由議論 \\
    \end{tabularx}
\end{table}%

\section*{5月14日 (木)}
\vspace{-10pt}
\begin{table}[H]
    \centering
    \begin{tabularx}{0.8\linewidth}{llY}
        \toprule
        \multicolumn{1}{l}{\textbf{時間}}
        &\multicolumn{1}{l}{\textbf{講演者}}
        &\multicolumn{1}{l}{\textbf{タイトル}} \\
        \midrule
        10:00-12:00 &\spkB &\titleB   \\ %\href{https://youtu.be/K9ZWQPkibxk}{\titleC} \\
        12:00-13:30 & & 昼休憩 \\
        13:30-15:00 &\spkC &\titleC  \\ %\href{https://youtu.be/rhHHrmUZW7s}{\titleC} \\
        15:30-17:00 &\spkC &\titleC  \\ %\href{https://youtu.be/xow1kzEf_n4}{\titleD} \\
        17:00-18:00 & &自由議論 \\
    \end{tabularx}
\end{table}%

\section*{5月15日 (金)}
\vspace{-10pt}
\begin{table}[H]
    \centering
    \begin{tabularx}{0.8\linewidth}{llY}
        \toprule
        \multicolumn{1}{l}{\textbf{時間}}
        &\multicolumn{1}{l}{\textbf{講演者}}
        &\multicolumn{1}{l}{\textbf{タイトル}} \\
        \midrule
        10:00-11:30 &\spkD &\titleD   \\ %\href{https://youtu.be/K9ZWQPkibxk}{\titleC} \\
        11:30-13:00 & & 昼休憩 \\
        13:00-14:30 &\spkD &\titleD  \\ %\href{https://youtu.be/rhHHrmUZW7s}{\titleC} \\
        14:30-16:30 & &自由議論  \\ %\href{https://youtu.be/xow1kzEf_n4}{\titleD} \\
    \end{tabularx}
\end{table}%

\section*{\underline{概要集}}

\section*{5月12日（火）}

% \TBA
\abst{\titleE}{\spkE}{\instE}{13:00-14:30, 15:00-16:30}{
    量子多体系の様々な性質は、そのエンタングルメント構造に反映されると考えられる。
    実際、エンタングルメント・エントロピーをはじめとする、エンタングルメントに関連した量を調べることにより、量子多体系の普遍的な性質を抽出できることが知られている。このような考え方は、トポロジカル・エンタングルメント・エントロピー等に代表されるように、量子多体系の理解に大きく貢献してきた。
    本講演では、multipartite entanglement、すなわち量子系を三つ以上の部分系に分割した際に現れる量子相関に焦点を当てる。特に、トポロジカル秩序をもつ基底状態について、multipartite entanglementを用いることで、どのような情報を抽出できるかを議論したい。    % \\\relax
}

\abst{\titleF}{\spkF}{\instF}{17:00-18:30}{
    本講演では、弦理論の摂動的定式化を与える世界面理論のトポロジカルな側面について議論します。弦理論の構成にはGSO射影やオリエンティフォールドといった操作が含まれます。これらは古くから知られていますが、近年の可逆な場の理論や量子異常に関する進展により、より体系的に理解できるようになりました。また、混成弦理論の構成に現れるカイラルな共形場理論を取り扱う上でも、量子異常やボゾン化・フェルミオン化に関する近年の発展が重要な役割を果たします。本講演では、特に10次元の弦理論に焦点を当て、II型弦理論のような時空の超対称性を保つ理論に加えて、0型弦理論や杉本弦理論のような超対称性を持たない理論についても議論します。世界面理論における量子異常や可逆な場の理論との結合の観点から、10次元の各種弦理論を導入し、標的空間に課されるトポロジカルな制約について概説します。時間が許せば、Dブレーン電荷のK理論による分類や、混成弦理論における非超対称ブレーンに関する話題にも触れる予定です。  
}


\newpage

\section*{5月13日（水）}

% \TBA
\abst{\titleF}{\spkF}{\instF}{10:00-11:30}{
    本講演では、弦理論の摂動的定式化を与える世界面理論のトポロジカルな側面について議論します。弦理論の構成にはGSO射影やオリエンティフォールドといった操作が含まれます。これらは古くから知られていますが、近年の可逆な場の理論や量子異常に関する進展により、より体系的に理解できるようになりました。また、混成弦理論の構成に現れるカイラルな共形場理論を取り扱う上でも、量子異常やボゾン化・フェルミオン化に関する近年の発展が重要な役割を果たします。本講演では、特に10次元の弦理論に焦点を当て、II型弦理論のような時空の超対称性を保つ理論に加えて、0型弦理論や杉本弦理論のような超対称性を持たない理論についても議論します。世界面理論における量子異常や可逆な場の理論との結合の観点から、10次元の各種弦理論を導入し、標的空間に課されるトポロジカルな制約について概説します。時間が許せば、Dブレーン電荷のK理論による分類や、混成弦理論における非超対称ブレーンに関する話題にも触れる予定です。
}

\abst{\titleA}{\spkA}{\instA}{13:00-14:30, 15:00-16:30}{
    \TBA
}

\newpage

\section*{5月14日（木）}

% \TBA
\abst{\titleB}{\spkB}{\instB}{10:00-12:00}{
    \TBA
}

\abst{\titleC}{\spkC}{\instC}{13:30-15:00, 15:30-17:00}{
    A型位相的弦理論の分配関数は，弦の世界面から標的空間への正則写像の “数え上げ” から
    定義される Gromov-Witten 不変量の生成母関数です．一方，M-理論の観点からは同じ分配関数が
    標的空間内の D-ブレインの BPS 状態の “数え上げ” から定義される Gopakumar-Vafa 不変量を用いて
    表すことができると予想されています．前半の講演では Gromov-Witten 不変量と Gopakumar-Vafa 
    不変量の関係を物理的な議論から説明します．後半は具体例として分配関数に対する定数写像の寄与と
    Resolved conifold に対する位相的弦理論の分配関数の計算を紹介します．定数写像の寄与が
    平面分割の数え上げという簡明な記述を持つことを見るとともに，時間が許せば
    Resolved conifold の分配関数と U(1) ゲージ理論の Nekrasov 分配関数との対応関係を説明します．
}

\newpage

\section*{5月15日（金）}

% \TBA
\abst{\titleD}{\spkD}{\instD}{10:30-12:00, 13:00-14:30}{
    Higher topos theory provides a unified framework for homotopy theory and modern geometry. An oo-topos is an oo-category that is supposed to be a collection of higher-categorical sheaves (or `stacks'). We begin by exposing the motivations and fundamentals of the theory from a categorical perspective. As an application, we discuss the intrinsic construction of the Bauer-Furuta invariant---a stable homotopy theoretic refinement of the Seiberg-Witten invariant.
}

\end{document}
